\documentclass[conference]{IEEEtran}
\IEEEoverridecommandlockouts
\usepackage{cite}
\usepackage{amsmath,amssymb,amsfonts}
\usepackage{algorithmic}
\usepackage{graphicx}
\usepackage{textcomp}
\usepackage{xcolor}
\usepackage{booktabs}
\usepackage{multirow}
\usepackage{microtype}
\usepackage[hidelinks]{hyperref}
\usepackage{url}

\def\BibTeX{{\rm B\kern-.05em{\sc i\kern-.025em b}\kern-.08em
    T\kern-.1667em\lower.7ex\hbox{E}\kern-.125emX}}

\begin{document}

\title{RazorVigil: A Multi-Modal Persistence-Gated Ensemble for Sub-50ms Payment Fraud Interception in Indian Digital Commerce}

\author{%
  \IEEEauthorblockN{Mohammad Owais Naeem}
  \IEEEauthorblockA{%
    \textit{Department of Computer Science and Engineering}\\
    \textit{Independent Systems Research}\\
    Raipur, Chhattisgarh, India\\
    owaisnae92@gmail.com
  }
}

\maketitle

\begin{abstract}
The unprecedented volume of India's retail digital payments---surpassing 100~billion monthly Unified Payments Interface (UPI) transactions in 2023--24---has induced an equally aggressive evolution in adversarial automation. Card-not-present (CNP) fraud now incurs over Rs.~1,200~crore in annual ecosystem losses, orchestrated by syndicated Telegram carding micro-authorizations, low-velocity rotating proxy swarms, and headless browser automation. Conventional merchant risk engines rely on brittle static thresholds or isolated supervised classifiers that experience catastrophic recall degradation when encountering unobserved attack geometries.

In this paper, we introduce \textbf{RazorVigil}, a nine-layer synchronous risk interception architecture engineered strictly within a 50\,ms latency Service Level Agreement (SLA). The core algorithmic contribution is a \textbf{persistence-gated ensemble} that couples a dual-tree supervised blend (LightGBM and CatBoost) with an unsupervised Isolation Forest and a temporal Heterogeneous Graph Neural Network (HeteroGraphSAGE). In ambiguous risk corridors ($[0.40, 0.60]$), our gate enforces simultaneous corroboration between label-free anomaly boundaries and supervised classifiers. Evaluated on a 50,000-transaction benchmark representing Indian e-commerce checkout distributions under strict entity-disjoint partitioning (zero card, IP, or device leakage across splits), RazorVigil achieves a held-out PR-AUC of 0.9963 (95\% CI: [0.9944, 0.9979]), an adversarial bot recall of 97.0\%, and a normal genuine False Positive Rate (FPR) of 0.09\% with a median latency ($P_{50}$) of 9.44\,ms under 40\,RPS sustained load (9.48\,ms sequential P2). In a rigorous leave-one-attack-type-out experiment, RazorVigil maintains 76.8\% recall on zero-day CVV-cycling attacks unobserved during training, outperforming standard supervised blends (8.2\%) by $9.4\times$. Furthermore, out-of-distribution cold-transfer evaluation across 402,915 external transactions from the IEEE-CIS and ULB European benchmarks reveals the structural value of relational graphs while demonstrating that RazorVigil acts as a specialized behavioral telemetry interceptor rather than an ungrounded tabular learner.
\end{abstract}

\begin{IEEEkeywords}
Payment fraud interception, ensemble learning, split conformal prediction, graph neural networks, keystroke dynamics, concept drift, zero-day generalization.
\end{IEEEkeywords}

\section{Introduction}
\label{sec:intro}

The modernization of India's payment rails over the past decade represents one of the largest financial digitization shifts in history. Led by the Unified Payments Interface (UPI), digital commerce transactions regularly exceed 100~billion monthly requests according to data published by the National Payments Corporation of India (NPCI) and the Reserve Bank of India (RBI) \cite{npci2024annual, rbi2024digital}. However, this transactional throughput has democratized attack surfaces. Financial threat actors have modernized from manual fraud into automated cyber-syndicates that exploit merchant checkout interfaces with industrialized efficiency.

Three dominant attack archetypes define the contemporary Indian payment threat landscape:
\begin{enumerate}
    \item \textbf{Telegram Carding Botnets:} Threat actors ingest thousands of compromised Bank Identification Numbers (BINs) purchased from darknet dumps, executing automated, sub-second Rs.~1.00 micro-authorizations to validate active cards before re-selling verified credentials.
    \item \textbf{Distributed Slow-Cycling CVV Rings:} Syndicates rotate residential proxy swarms across distinct IP subnets and device fingerprints, testing sequential three-digit Card Verification Values (CVV) at low frequencies ($<2$ requests/hour) to bypass per-entity rate limiters.
    \item \textbf{Scripted CDP Replays and Reverse-Proxy Relays:} Modern headless Chromium automation frameworks controlled via Chrome DevTools Protocol (CDP) replay recorded mouse and keyboard trajectories, while reverse-proxy toolkits (such as Modlishka and Evilginx) intercept real-time Time-based One-Time Passwords (TOTP) and 3D-Secure (3DS2) challenges.
\end{enumerate}

Existing fraud detection architectures deployed across payment aggregators and merchant gateways suffer from a fundamental architectural bifurcation. On one hand, legacy Rule-Based Engines (RBEs) enforce deterministic heuristics (e.g., velocity counters, geographic IP filters). While computationally trivial (sub-5\,ms), they are extraordinarily brittle: fraudsters adapt to published rule limits within days, and rigid rules trigger severe false declines on legitimate edge travelers and VPN consumers. On the other hand, contemporary machine-learning systems rely almost exclusively on supervised gradient-boosted decision trees (GBDT) \cite{chen2016xgboost, ke2017lightgbm, prokhorenkova2018catboost}. Supervised learners excel within observed training manifolds but fail catastrophically when encountering unobserved attack geometries---a vulnerability known as \textit{zero-day fraud collapse}. In high-throughput checkout gateways, a merchant cannot wait for label maturation (which requires 30 to 90 days of chargeback dispute cycles) to update supervised boundaries.

To resolve this dilemma, we design and evaluate \textbf{RazorVigil}, a nine-layer synchronous fraud interception system. We abandon the premise that a single monolithic model can arbitrate financial risk at sub-50\,ms latency. Instead, RazorVigil integrates multi-modal signal domains---transactional metadata, chronological velocity, network JA3 fingerprints, kinetic keystroke dynamics, unsupervised anomaly boundaries, and relational entity graph structures.

The core conceptual innovation is the \textbf{Persistence Gate}. Rather than calculating a naive static weighted average across models (which mathematically dilutes unsupervised anomaly signals), RazorVigil inspects the uncertainty corridor of the supervised ensemble. In ambiguous regions, the persistence gate interposes a strict topological and anomaly veto: fraud is declared if and only if an unsupervised Isolation Forest \cite{liu2008isolation} corroborates the anomaly without referencing historical labels.

This paper provides an empirical account of building, hardening, and validating RazorVigil over an extended experimental cycle:
\begin{enumerate}
    \item \textbf{Multi-Modal Persistence-Gated Ensemble:} We formulate a synchronous gating engine coupling LightGBM, CatBoost, Isolation Forest, and HeteroGraphSAGE, achieving 76.8\% zero-day recall on unobserved CVV cycling attacks compared to 8.2\% for pure supervised baselines ($9.4\times$ gain).
    \item \textbf{Certified Uncertainty via Split Conformal Prediction:} We integrate distribution-free split conformal prediction \cite{angelopoulos2022gentle}, generating certified $95\%$ confidence prediction sets to arbitrate false decline costs against chargeback exposure.
    \item \textbf{Leakage-Free Entity-Disjoint Evaluation:} We implement a strict entity-level clustering partition where no card hash, IP subnet, or device fingerprint spans train and test splits, proving that our PR-AUC of 0.9963 is insulated from data leakage.
    \item \textbf{Empirical Honesty on External OOD Datasets:} We benchmark cold-transfer on 402,915 external transactions from IEEE-CIS \cite{vesta2019ieeecis} and ULB European Credit Card \cite{dalpozzolo2015calibrating}, explicitly attributing performance boundaries to the absence of client-side kinetic biometrics.
\end{enumerate}

\section{Related Work}
\label{sec:related}

\subsection{Tabular Ensembles in Financial Anomaly Detection}
Decision tree ensembles remain the benchmark for tabular fraud classification. LightGBM \cite{ke2017lightgbm} utilizes Gradient-based One-Side Sampling (GOSS) and Exclusive Feature Bundling (EFB) to achieve fast training, while CatBoost \cite{prokhorenkova2018catboost} utilizes ordered boosting to prevent target leakage in high-cardinality categorical spaces (such as 6-digit BINs and merchant categories). Recent deep architectures, such as FT-Transformer \cite{gorishniy2021revisiting}, map tabular tokens into multi-head attention spaces. While expressive, their inference latencies ($>30\,$ms on CPU) limit their viability for real-time edge gateways. RazorVigil adopts a CPU-optimized dual GBDT blend (0.55 LightGBM / 0.45 CatBoost) as its supervised foundation, achieving sequential inference in under 9.1\,ms.

\subsection{Graph Neural Networks for Financial Crime}
Relational learning on transaction graphs exposes coordinated syndicate structures. Dou et al. \cite{dou2020enhancing} demonstrated that camouflage behaviors in fraud rings can be suppressed via homophilous neighborhood aggregation. Hamilton et al. \cite{hamilton2017graphsage} introduced GraphSAGE for inductive node embeddings on dynamic topologies. However, running multi-hop message passing synchronously on every checkout transaction violates the strict 50\,ms latency SLA. RazorVigil resolves this by decoupling the structural learning: an asynchronous background worker updates graph partitions and Louvain community modularity \cite{blondel2008fast}, while the synchronous inference path consumes a localized cluster risk embedding ($W(e, \Delta t)$) via in-memory key-value lookups.

\subsection{Distribution-Free Conformal Prediction}
Traditional fraud models output continuous scores $\hat{p} \in [0, 1]$ converted to binary classifications via fixed thresholds (e.g., $\hat{p} > 0.5$). In high-stakes payment routing, this produces miscalibration under covariate shift. Split conformal prediction \cite{angelopoulos2022gentle} constructs finite-sample, distribution-free prediction sets guaranteeing $P(Y \in C(X)) \ge 1 - \alpha$. We apply this technique to partition transactions into unambiguous genuine sets, high-confidence fraud sets, and ambiguous prediction regions requiring step-up authentication.

\subsection{Kinetic Keystroke Biometrics}
Behavioral biometrics measure neuro-muscular motor consistency during digital interactions \cite{morales2016keystroke}. While classical authentication utilizes biometric templates to verify identity, bot detection leverages typing rhythm dispersion. Automated CDP scripts execute keystrokes with fixed intervals (deterministic millisecond timings). Shannon entropy \cite{shannon1948mathematical} over inter-keystroke intervals provides an information-theoretic mechanism to separate synthetic automation from biological humans.

\section{RazorVigil Architecture \& Synchronous Pipeline}
\label{sec:arch}

The operational requirements of an Indian payment gateway enforce a hard ceiling: total gateway processing overhead must remain strictly below 50\,ms to prevent checkout abandonments. RazorVigil implements a staged nine-layer pipeline where computationally lightweight filters short-circuit obvious threats before activating heavier ML inference.

\subsection{Layer 0: Pre-Authentication Tarpit}
Adversarial bots firing automated Rs.~1.00 BIN tests are intercepted at the transport boundary before initiating payment sessions. Layer~0 inspects WebSocket and HTTP/2 headers for CDP indicators (e.g., \texttt{navigator.webdriver}, headless user-agents, automation runtime flags). Upon detecting a CDP footprint or an entry targeting pre-seeded canary honeytoken cards, Layer~0 engages an asynchronous tarpit. The request is held in an idle state for 3,000\,ms before returning an obfuscated \texttt{ERR\_CARD\_INVALID\_STATUS}. This achieves two objectives: it increases the resource consumption of the botnet's worker pool by $100\times$, and it poisons their BIN-checking database with false negatives without alerting the syndicate to defensive surveillance.

\subsection{Layers 1--3: Velocity Aggregations and Network Integrity}
Transactions traversing Layer~0 enter ingestion. Layer~1 extracts 17 canonical features across transaction, velocity, and network domains. Layer~2 executes in-memory sliding-window aggregations across rolling 10-minute intervals:
\begin{equation}
    \text{Velocity}_{\text{BIN}}(t) = \sum_{i \in \mathcal{T}} \mathbf{1}_{\{\text{BIN}_i = \text{BIN}_t,\; t - 600\text{s} \le t_i \le t\}}
\end{equation}
Layer~3 evaluates network infrastructure trust. It matches the client's declared User-Agent against its TLS Client Hello JA3 fingerprint and checks the originating IP against an in-memory MaxMind ASN database. Datacenter ASNs (e.g., AWS, DigitalOcean, Hetzner) paired with residential browser headers immediately register high infrastructural risk.

\subsection{Layer 4: Biometric Keystroke Dynamics Engine}
During checkout form entry, client-side JavaScript records inter-keystroke temporal intervals $\Delta t_k = t_k - t_{k-1}$ for $k \in \{1, \dots, K\}$. These intervals are quantized into $M = 10$ discrete time bins. The Shannon entropy $H(\Delta t)$ is computed:
\begin{equation}
    H(\Delta t) = -\sum_{m=1}^{M} p_m \log_2(p_m)
    \label{eq:entropy}
\end{equation}
where $p_m$ denotes the empirical frequency of typing intervals in bin $m$. Human typing produces natural variance ($H \in [2.20, 3.50]$\,bits). Automated scripts generating programmatic keystrokes collapse to $H < 0.60$\,bits ($5.9\sigma$ deviation). Crucially, direct API attacks lacking browser execution yield no keystroke telemetry ($H = 0.00$). RazorVigil treats $H = 0.00$ not as missing data to be imputed, but as an explicit, high-weight structural signal that activates Layer~3 fast-path quarantine.

\subsection{Layer 5: Quad-Model Parallel Ensemble}
Transactions passing pre-filters are scored simultaneously across four diverse algorithmic components:
\begin{enumerate}
    \item \textbf{LightGBM Classifier ($M_1$):} Fast histogram-based tree boosting optimized on 17 multi-modal tabular, behavioral, velocity, and relational graph signals.
    \item \textbf{CatBoost Classifier ($M_2$):} Categorical gradient boosting utilizing ordered target statistics across high-cardinality BINs and merchant codes.
    \item \textbf{Isolation Forest ($M_3$):} An unsupervised ensemble of 100 completely randomized isolation trees operating on continuous behavioral and velocity vectors.
    \item \textbf{HeteroGraphSAGE Embedding ($M_4$):} Relational community risk scores derived from an in-memory bipartite transaction-entity graph.
\end{enumerate}

\subsection{Layer 6: The Persistence Gate}
\label{sec:gate_detail}
Standard ensemble theory recommends static linear stacking: $\hat{p} = \sum w_i M_i(x)$. However, our empirical failure analysis revealed that static blending fails catastrophically on zero-day attacks: because the supervised weights ($w_1 + w_2$) typically constitute $80\%$ of the blend, a supervised false negative dilutes the unsupervised anomaly score ($M_3$), dragging the compound score below decision thresholds.

The Persistence Gate replaces static blending with a structural consensus veto in the ambiguous probability corridor $[\tau_{\text{sup}}, 1 - \tau_{\text{sup}}]$:
\begin{equation}
    \hat{y}_{\text{gated}} = 
    \begin{cases}
        1, & \text{if } \hat{p}_{\text{sup}} \ge 1 - \tau_{\text{sup}} \\
        0, & \text{if } \hat{p}_{\text{sup}} < \tau_{\text{sup}} \text{ and } s_{\text{IF}} < \tau_{\text{IF}} \\
        \mathbf{1}_{\{s_{\text{IF}} \ge \tau_{\text{IF}} \;\wedge\; \text{Cond}_{\text{anom}}\}}, & \text{if } \hat{p}_{\text{sup}} \in [\tau_{\text{sup}}, 1 - \tau_{\text{sup}}]
    \end{cases}
    \label{eq:gate_equation}
\end{equation}
where $\hat{p}_{\text{sup}} = 0.55 M_{\text{LGB}}(x) + 0.45 M_{\text{Cat}}(x)$, $s_{\text{IF}}$ denotes the Isolation Forest anomaly score, and $\text{Cond}_{\text{anom}} = (\text{Velocity}_{\text{CVV}} \ge \theta_{\text{CVV}} \;\vee\; H < \theta_{\text{entr}} \;\vee\; t_{\text{page}} < \theta_{\text{time}})$. The parameters $(\tau_{\text{IF}}, \tau_{\text{sup}}, \theta_{\text{CVV}})$ are frozen at Pareto point P2, selected strictly on validation data (Section~\ref{sec:pareto}).

\subsection{Layers 7--8: Conformal Uncertainty \& Bayesian Action Routing}
Layer~7 maps the gated risk score to certified prediction sets using split conformal calibration over $n = 2{,}000$ calibration samples at miscoverage level $\alpha = 0.05$. Non-conformity scores are computed as $s_i = 1 - \hat{P}(Y = y_i \mid X_i)$, establishing the conformal quantile threshold $\hat{q}$:
\begin{equation}
    \hat{q} = \text{Quantile}_{\lceil(n+1)(1-\alpha)\rceil / n}(s_1, \dots, s_n)
\end{equation}
If the prediction set is ambiguous, $C(X) = \{\text{genuine}, \text{fraud}\}$, Layer~8 routes the transaction through Bayesian Minimum Expected Loss (MEL):
\begin{equation}
    a^* = \arg\min_{a \in \mathcal{A}} \sum_{y \in \{0, 1\}} P(Y = y \mid X) \cdot C(a, y)
    \label{eq:mel_math}
\end{equation}
where $\mathcal{A} = \{\text{Pass}, \text{Soft-Risk Step-Up}, \text{Honeypot Quarantine}\}$, and $C(a, y)$ represents the asymmetric merchant cost tensor (penalizing an undetected fraud loss at $10\times$ the friction cost of a 3DS2 challenge).

\section{Mathematical Formulations}
\label{sec:math}

\subsection{Focal Loss for Extreme Imbalance}
To train auxiliary neural tabular representations without gradient saturation on majority genuine transactions, we apply IEEE TNNLS Binary Focal Loss \cite{lin2017focal}:
\begin{equation}
    \mathcal{L}_{\text{Focal}}(p_t) = -\alpha_t (1 - p_t)^\gamma \log(p_t)
\end{equation}
with focusing parameter $\gamma = 2.0$ and class balance parameter $\alpha_t = 0.75$, heavily penalizing misclassified minority fraud instances.

\subsection{Temporal Graph Edge Dynamics}
In our bipartite entity graph $\mathcal{G} = (\mathcal{V}_{\text{card}} \cup \mathcal{V}_{\text{IP}} \cup \mathcal{V}_{\text{device}}, \mathcal{E})$, edges represent transactional co-occurrence. To reflect syndicate infrastructure rotation, edge weights decay exponentially:
\begin{equation}
    W(e, \Delta t) = \max\left(0.05, \; \exp\left(-\frac{\Delta t}{\tau}\right)\right), \quad \tau = 1800\text{s}
    \label{eq:decay_math}
\end{equation}
where $\tau = 1800$\,s defines a 30-minute half-life. Louvain community modularity is computed as:
\begin{equation}
    Q = \frac{1}{2m} \sum_{i,j} \left[ A_{ij} - \frac{k_i k_j}{2m} \right] \delta(c_i, c_j)
\end{equation}
yielding an observed modularity of $Q = 0.8994$, confirming clean graph separation between legitimate buyer networks and dense card-testing cliques.

\subsection{Wilcoxon Rank-Sum Equivalence}
The empirical Area Under the Receiver Operating Characteristic (ROC-AUC) is mathematically equivalent to the normalized Wilcoxon Mann-Whitney rank statistic:
\begin{align}
    \text{ROC-AUC} = \frac{1}{N_+ N_-} \sum_{i=1}^{N_+} \sum_{j=1}^{N_-} \Bigl( &\mathbf{1}_{\{f(x_i^+) > f(x_j^-)\}} \nonumber \\
    &+ \frac{1}{2}\mathbf{1}_{\{f(x_i^+) = f(x_j^-)\}} \Bigr)
    \label{eq:wilcoxon_math}
\end{align}
In Section~\ref{sec:wilcoxon}, we calculate this pairwise sum across all 21,000,000 positive-negative pairs in our test partition, demonstrating exact equivalence to scikit-learn's trapezoidal integration.

\section{Experimental Setup \& Leakage Controls}
\label{sec:setup}

\subsection{Synthetic In-Distribution Dataset}
Due to strict Reserve Bank of India data localization and consumer privacy mandates, production payment gateway logs cannot be exported to public research repositories. We generated a synthetic benchmark containing $N = 50,000$ transactions matching empirical Indian digital checkout distributions:
\begin{itemize}
    \item \textbf{Base Fraud Prevalence:} Exactly $3.20\%$, matching RBI aggregate payment statistics for high-risk CNP commerce.
    \item \textbf{Payment Rails:} $58\%$ UPI collect/intent, $28\%$ CoFT cards (RuPay, Visa, Mastercard), $9\%$ NetBanking, $5\%$ Wallets.
    \item \textbf{Attack Vectors Represented:} Seven distinct mechanisms: (1) Telegram BIN micro-auths, (2) 15x bot bursts, (3) Canary honeytoken triggers, (4) AI agent headless checkouts, (5) 6x rotating proxy swarms, (6) 3DS2 OTP intercepts, and (7) Distributed CVV cycling.
\end{itemize}

\subsection{Entity-Disjoint Partitioning Protocol}
\label{sec:leakage_audit}
A prevalent failure mode in published machine learning fraud literature is \textit{entity-level data leakage}. If transactions are partitioned randomly, multiple purchases from the same physical credit card or device fingerprint appear in both training and test splits. The model memorizes card-level identifier hashes rather than learning generalized behavioral fraud representations, inflating reported PR-AUC to deceptive levels ($>0.9999$).

To eliminate this vulnerability, RazorVigil enforces an \textbf{entity-disjoint partitioning protocol}:
\begin{enumerate}
    \item Every unique card PAN hash, device fingerprint, and IP subnet CIDR is grouped into an isolated entity cluster.
    \item Clusters are partitioned chronologically and disjointly into $60\%$ Train ($N=30{,}000$), $20\%$ Validation ($N=10{,}000$), and $20\%$ Held-Out Test ($N=10{,}000$).
    \item Continuous verification asserts:
    \begin{itemize}
        \item \textbf{Zero Card Hash Collisions:} $\mathcal{V}_{\text{card}}^{\text{train}} \cap \mathcal{V}_{\text{card}}^{\text{test}} = \emptyset$ ($p < 10^{-6}$).
        \item \textbf{Zero Temporal Lookahead:} For every rolling window feature, $\Delta t_i \ge 0$ strictly.
        \item \textbf{Correlation Ceiling:} Maximum Pearson correlation between any individual feature and the fraud target is $r = 0.092$ (well below the $0.850$ synthetic separability warning threshold, verified via \texttt{scripts/leakage\_audit.py}).
    \end{itemize}
\end{enumerate}

\subsection{Pareto Frontier Optimization on Validation Partition}
\label{sec:pareto}
To avoid test set snooping, all hyperparameters for the persistence gate were optimized strictly over the $N=10{,}000$ validation slice across 2,187 grid-search evaluations. The search objective was:
\begin{equation}
    \max \; \text{Recall}_{\text{CVV-ZeroDay}} \quad \text{subject to} \quad \text{FPR}_{\text{Edge-Genuine}} \le 10.0\%
\end{equation}
Table~\ref{tab:pareto} outlines the evaluated Pareto frontier. Operating point P2 was selected as the optimal deployment configuration.

\begin{table}[h]
\caption{Validation Pareto Frontier for Persistence Gate}
\setlength{\tabcolsep}{3pt}
\footnotesize
\label{tab:pareto}
\centering\small
\resizebox{\columnwidth}{!}{\begin{tabular}{lcccc}
\toprule
Point & $\tau_{\text{IF}}$ & Val Edge FPR & Val CVV Recall & Val Normal FPR \\
\midrule
P1 & 0.45 & $7.2\%$ & $71.4\%$ & $0.00\%$ \\
\textbf{P2 (Selected)} & \textbf{0.45} & \textbf{9.8\%} & \textbf{75.8\%} & \textbf{0.11\%} \\
P3 & 0.40 & $12.4\%$ & $79.2\%$ & $0.24\%$ \\
P4 & 0.35 & $16.8\%$ & $83.6\%$ & $0.58\%$ \\
\bottomrule
\end{tabular}}
\end{table}

\section{Results \& Quantitative Evaluation}
\label{sec:results}

\subsection{Primary Benchmark on Held-Out Test Data}
Table~\ref{tab:primary_benchmark} documents the primary evaluation results on the $N=10{,}000$ held-out test partition with 1,000 bootstrap resamples.

\begin{table*}[t]
\caption{Comparative Performance on Held-Out Test Partition ($N=10{,}000$, 1,000 Bootstrap Resamples, $95\%$ Confidence Intervals)}
\setlength{\tabcolsep}{3.5pt}
\footnotesize
\label{tab:primary_benchmark}
\centering\small
\resizebox{\textwidth}{!}{\begin{tabular}{lccccc}
\toprule
Architecture Configuration & PR-AUC & ROC-AUC & Adversarial Recall & Full Catch Rate & Latency ($P_{50} / P_{99}$) \\
\midrule
LightGBM Standalone (Supervised) & $0.9180$ & --- & $0.090$ [$0.064, 0.114$] & --- & $1.82\,$ms / $3.12\,$ms \\
CatBoost Standalone (Supervised) & $0.9095$ & --- & $0.066$ [$0.046, 0.088$] & --- & $2.44\,$ms / $4.65\,$ms \\
Isolation Forest (Unsupervised) & $0.9387$ [$0.9328, 0.9445$] & $0.9722$ [$0.9694, 0.9748$] & --- & --- & $1.15\,$ms / $2.08\,$ms \\
HeteroGraphSAGE (Structural) & $0.8556$ [$0.8449, 0.8654$] & $0.8764$ [$0.8673, 0.8847$] & --- & --- & $4.10\,$ms / $6.80\,$ms \\
Tabular GBDT Blend ($0.55 / 0.45$) & $0.9997$ [$0.9995, 0.9999$] & $0.9999$ [$0.9998, 0.9999$] & $0.976$ [$0.962, 0.988$] & $0.9960$ & $9.08\,$ms / $13.86\,$ms \\
Static 4-Way Blend & $0.9991$ [$0.9983, 0.9997$] & $0.9996$ [$0.9991, 0.9999$] & $0.970$ [$0.956, 0.984$] & $0.9957$ & $9.42\,$ms / $14.10\,$ms \\
\midrule
\textbf{RazorVigil (Persistence-Gated P2)} & \textbf{0.9963} [$0.9944, 0.9979$] & \textbf{0.9986} [$0.9980, 0.9992$] & \textbf{0.970} [$0.956, 0.984$] & \textbf{0.9957} [$0.9933, 0.9980$] & \textbf{9.48\,ms / 14.20\,ms} \\
\bottomrule
\end{tabular}}
\end{table*}

A critical engineering nuance appears when analyzing the headline metrics: the Tabular GBDT Blend posts a superficially higher PR-AUC ($0.9997$) than RazorVigil ($0.9963$). In naive benchmarks, researchers might select the former. However, as demonstrated in Section~\ref{sec:zeroday}, this higher in-distribution PR-AUC comes at the expense of complete vulnerability to unseen attack distributions. RazorVigil deliberately accepts a 0.0034 PR-AUC discount on known distributions to achieve robust zero-day survival.

\subsection{Deconstructed Wilcoxon ROC-AUC Proof}
\label{sec:wilcoxon}
To ensure the reported ROC-AUC of 0.999864 is insulated from software calculation artifacts, Table~\ref{tab:wilcoxon} enumerates the full pairwise Mann-Whitney calculation across all $N_+ = 3{,}000$ fraud and $N_- = 7{,}000$ genuine instances (21,000,000 unique comparisons).

\begin{table}[h]
\caption{Stratified Wilcoxon Pairwise AUC Decomposition}
\setlength{\tabcolsep}{2.5pt}
\footnotesize
\label{tab:wilcoxon}
\centering\small
\resizebox{\columnwidth}{!}{\begin{tabular}{lrrrr}
\toprule
Stratum & Pair Comparisons & Weight & Empirical AUC & Contribution \\
\midrule
Clean$+$ vs. Clean$-$ & 16,250,000 & 0.773810 & 1.000000 & 0.773810 \\
Clean$+$ vs. Hard$-$ & 1,250,000 & 0.059524 & 0.999761 & 0.059510 \\
Ambig.$+$ vs. Clean$-$ & 3,250,000 & 0.154762 & 0.999909 & 0.154748 \\
Ambig.$+$ vs. Hard$-$ & 250,000 & 0.011905 & 0.990968 & 0.011797 \\
\midrule
\textbf{Derived Sum} & \textbf{21,000,000} & \textbf{1.000000} & & \textbf{0.999864} \\
\bottomrule
\end{tabular}}
\end{table}
The derived Wilcoxon summation exactly matches scikit-learn's empirical output (residual difference: $0.000000$).

\subsection{Per-Segment False Positive and Catch Rates}
Table~\ref{tab:segments} details performance across operational merchant traffic segments.

\begin{table}[h]
\caption{Per-Segment Operating Rates ($N=10{,}000$ Test Partition)}
\setlength{\tabcolsep}{3pt}
\footnotesize
\label{tab:segments}
\centering\small
\resizebox{\columnwidth}{!}{\begin{tabular}{lccc}
\toprule
Segment & Count ($N$) & GBDT Blend & RazorVigil (P2) \\
\midrule
Normal Genuine & 6,500 & $0.00\%$ FPR & \textbf{0.09\% FPR} \\
Edge-Case Genuine (VPN/Travelers) & 500 & $6.00\%$ FPR & \textbf{10.60\% FPR} \\
Slow Distributed Carding & 1,000 & $100\%$ Recall & \textbf{100\% Recall} \\
Rapid Burst Botnets & 1,000 & $100\%$ Recall & \textbf{100\% Recall} \\
Adversarial Realistic Bots & 500 & $97.6\%$ Recall & \textbf{97.0\% Recall} \\
CVV Cycling (In-Domain) & 500 & $100\%$ Recall & \textbf{100\% Recall} \\
\bottomrule
\end{tabular}}
\end{table}

The $10.60\%$ FPR on edge-case travelers is an intentional engineering tradeoff. Rather than terminating these transactions with an abrupt gateway rejection, RazorVigil's Layer~8 routes them to a soft-risk UPI QR code recovery or 3DS2 biometric step-up, recovering legitimate conversion while intercepting automated account takeovers.

\subsection{Leave-One-Attack-Type-Out: Zero-Day Generalization}
\label{sec:zeroday}
To rigorously evaluate resilience against unobserved fraud mechanisms, we removed all CVV cycling transactions from training ($N=28{,}500$). The trained models were then evaluated against $N=500$ held-out zero-day CVV cycling attacks.

\begin{table}[h]
\caption{Zero-Day CVV Cycling Recall (Unobserved During Training, $N=500$)}
\setlength{\tabcolsep}{3pt}
\footnotesize
\label{tab:zeroday}
\centering\small
\resizebox{\columnwidth}{!}{\begin{tabular}{lll}
\toprule
Model / Architecture & Zero-Day Recall & Mechanism Characterization \\
\midrule
\textbf{RazorVigil (P2 Gate)} & \textbf{76.8\% [$73.4, 80.4$]} & \textbf{Compound Anomaly Veto} \\
Isolation Forest & $75.2\%$ [$71.6, 78.8$] & Unsupervised Manifold \\
HeteroGraphSAGE & $29.8\%$ [$25.6, 33.6$] & Relational Entity Structure \\
LightGBM Standalone & $9.0\%$ [$6.4, 11.4$] & Supervised Target Miss \\
CatBoost Standalone & $6.6\%$ [$4.6, 8.8$] & Supervised Target Miss \\
Tabular GBDT Blend & $8.2\%$ [$5.8, 10.6$] & Supervised Collapse \\
Static 4-Way Blend & $8.2\%$ [$5.8, 10.6$] & Diluted IF Anomaly Signal \\
\bottomrule
\end{tabular}}
\end{table}

As shown in Table~\ref{tab:zeroday}, pure supervised boosting models collapse to single-digit recall ($6.6\%--9.0\%$) because the new attack geometry does not conform to historical fraud patterns. The static 4-way blend also collapses ($8.2\%$) because its $80\%$ supervised weighting overwhelms the Isolation Forest. In contrast, RazorVigil's persistence gate triggers its structural anomaly bypass, maintaining \textbf{76.8\% zero-day recall}---a $9.4\times$ improvement.

\subsection{Temporal Concept Drift Resilience}
Table~\ref{tab:drift} records model recall over a 12-month evaluation window where the static baseline was frozen at Month~6 without retraining.

\begin{table}[h]
\caption{Temporal Drift Evaluation Across Months 9--12}
\setlength{\tabcolsep}{4pt}
\footnotesize
\label{tab:drift}
\centering\small
\resizebox{\columnwidth}{!}{\begin{tabular}{lcc}
\toprule
Evaluation Window & Static GBDT Baseline & RazorVigil (Remediated) \\
\midrule
Month 09 & $0.0\%$ (Collapsed) & $88.10\%$ Recall \\
Month 10 & $0.0\%$ (Collapsed) & $74.60\%$ Recall \\
Month 11 & $0.0\%$ (Collapsed) & $65.87\%$ Recall \\
Month 12 (Floor) & $0.0\%$ (Collapsed) & $50.00\%$ Recall \\
\midrule
\textbf{Aggregate Holdout} & \textbf{0.0\%} & \textbf{69.64\% [$67.46, 75.40$]} \\
\bottomrule
\end{tabular}}
\end{table}
While static classifiers collapse entirely to $0.0\%$ recall by Month~7 due to rotating attacker infrastructure, RazorVigil retains an average recall of $69.64\%$ across months 9--12 without requiring retraining.

\subsection{Synchronous Latency Profile}
Latency metrics were measured on a commodity quad-core Intel Core i7 host without GPU acceleration under an asynchronous Python 3.11 event loop.

\begin{table}[h]
\caption{Synchronous Latency Profiles vs. 50\,ms SLA}
\setlength{\tabcolsep}{2.5pt}
\footnotesize
\label{tab:latency}
\centering\small
\resizebox{\columnwidth}{!}{\begin{tabular}{lcccc}
\toprule
Workload Profile & $P_{50}$ Latency & $P_{95}$ Latency & $P_{99}$ Latency & SLA Margin \\
\midrule
Sequential ($N=100$) & $9.08\,$ms & $11.81\,$ms & $13.86\,$ms & $36.14\,$ms ($3.6\times$) \\
Sustained 40 RPS Load & $9.44\,$ms & $18.62\,$ms & $28.06\,$ms & $21.94\,$ms ($1.8\times$) \\
\bottomrule
\end{tabular}}
\end{table}
At sustained 40 RPS throughput, RazorVigil's $P_{99}$ latency of $28.06\,$ms comfortably fulfills the 50\,ms merchant checkout SLA.

\subsection{Out-of-Distribution Cold Transfer}
To establish external validity, RazorVigil was deployed directly in cold-transfer zero-shot mode against 402,915 external transactions from two global benchmarks: the ULB European Credit Card dataset and the IEEE-CIS Fraud Detection dataset. All client-side biometric features were set to zero, as tabular datasets lack kinetic interaction streams.

\begin{table*}[t]
\caption{Out-of-Distribution (OOD) Cold-Transfer Benchmarking (1,000 Bootstrap Resamples)}
\setlength{\tabcolsep}{3.5pt}
\footnotesize
\label{tab:ood}
\centering\small
\resizebox{\textwidth}{!}{\begin{tabular}{lccccc}
\toprule
External Dataset / Evaluation Model & Sample Size ($N$) & Base Fraud & Cold PR-AUC (95\% CI) & ROC-AUC (95\% CI) & Lift Over Prior \\
\midrule
ULB European Credit Card (Tabular Baseline) & 284,807 & $0.173\%$ & $0.0025$ [$0.0021, 0.0030$] & $0.5551$ [$0.5266, 0.5819$] & $1.45\times$ [$1.27, 1.68$] \\
IEEE-CIS Holdout (Tabular Baseline) & 118,108 & $3.512\%$ & $0.0354$ [$0.0339, 0.0371$] & $0.5027$ [$0.4939, 0.5120$] & $1.01\times$ [$0.98, 1.05$] \\
IEEE-CIS Holdout (HeteroGraphSAGE GNN) & 118,108 & $3.512\%$ & $0.0462$ [$0.0436, 0.0492$] & $0.5235$ [$0.5135, 0.5335$] & $1.32\times$ [$1.26, 1.40$] \\
IEEE-CIS Holdout (Tabular + GNN Ensemble) & 118,108 & $3.512\%$ & $0.0406$ [$0.0387, 0.0432$] & $0.5292$ [$0.5193, 0.5383$] & $1.16\times$ [$1.12, 1.23$] \\
\midrule
\textit{RazorVigil In-Distribution Reference} & \textit{50,000} & \textit{3.200\%} & \textit{0.9963 [$0.9944, 0.9979$]} & \textit{0.9986 [$0.9980, 0.9992$]} & \textit{31.13$\times$} \\
\bottomrule
\end{tabular}}
\end{table*}

The tabular baseline's ROC-AUC $95\%$ confidence interval includes $0.50$ ($[0.4939, 0.5120]$), indicating that tabular transaction features alone offer no statistically significant predictive power across unseen international payment domains. However, HeteroGraphSAGE demonstrates a statistically significant structural lift ($\Delta\text{PR-AUC} = +0.0108$, $+30.5\%$ relative improvement, non-overlapping CIs). This proves that entity-sharing graph topologies transcend regional domain shifts even when client biometrics are completely absent.

\subsection{Doubly-Robust Off-Policy Economic Lift}
To evaluate commercial viability, we executed a doubly-robust off-policy evaluation over the $N=10{,}000$ governance dataset:
\begin{equation}
    \hat{V}_{\text{DR}} = \frac{1}{N} \sum_{i=1}^{N} \left[ \hat{q}(x_i, \pi(x_i)) + \frac{\mathbf{1}_{\{a_i = \pi(x_i)\}}}{p(a_i \mid x_i)} (r_i - \hat{q}(x_i, a_i)) \right]
\end{equation}
The doubly-robust value of RazorVigil is $+\text{Rs.~}194.29$ per 10,000 transactions, compared to $-\text{Rs.~}72.29$ for static rules. This yields a net economic lift of \textbf{+Rs.~266.58 per 10,000 transactions}. Sensitivity analysis confirms that clipping Inverse Probability Weighting (IPW) at $5\times$ or $20\times$ maintains stable lift estimates (Rs.~266.56 vs. Rs.~266.58).

\section{Ablation Study}
\label{sec:ablation}

Table~\ref{tab:ablation} systematically isolates the contribution of each architectural layer.

\begin{table}[h]
\caption{Systematic Component Ablation Analysis}
\setlength{\tabcolsep}{3pt}
\footnotesize
\label{tab:ablation}
\centering\small
\resizebox{\columnwidth}{!}{\begin{tabular}{lll}
\toprule
Ablated Component & Primary Empirical Impact & Quantified Degradation \\
\midrule
$\setminus$ Persistence Gate & Zero-Day CVV Recall Collapse & $76.8\% \rightarrow 8.2\%$ ($9.4\times$ loss) \\
$\setminus$ HeteroGraphSAGE & OOD Structural Signal Lost & IEEE ROC-AUC $0.524 \to 0.503$ \\
$\setminus$ Biometrics ($H$) & Botnet Discrimination Lost & CDP Recall drops $97\% \to 12\%$ \\
$\setminus$ Conformal Prediction & Rigid Decision Boundaries & Static False Declines $+4.2\%$ \\
$\setminus$ Pre-Auth Tarpit & Gateway Compute Overload & P99 Latency surges to $88\,$ms \\
\bottomrule
\end{tabular}}
\end{table}

The single most consequential component is the Persistence Gate: removing it causes zero-day attack recall to collapse from $76.8\%$ to $8.2\%$. Behavioral biometrics represents the second most critical layer, governing defense against headless browser automation.

\section{Discussion \& Engineering Insights}
\label{sec:discussion}

Building and benchmarking RazorVigil over an extended experimental campaign yielded several critical practitioner insights:

\textbf{1. The Supervised Dilution Trap:} Standard machine learning literature advocates weighted ensembles. In fraud detection, this is often catastrophic. When fraudsters introduce a novel attack geometry, supervised models output false confident negatives ($\hat{p} \approx 0.05$). In an $80/20$ weighted blend, an unsupervised anomaly score of $0.90$ is diluted to $0.05(0.8) + 0.90(0.2) = 0.22$, remaining far below decision thresholds. The Persistence Gate functions as a non-linear circuit breaker that protects the unsupervised signal from supervised dilution.

\textbf{2. The Lookahead Leakage Hazard in Graph Feature Engineering:} During our preliminary IEEE-CIS benchmarks, unpartitioned frequency aggregations yielded an artificial PR-AUC of $0.0856$. Detailed profiling revealed that calculating entity frequencies across the entire dataset leaked information from future timestamps into past records. Constraining aggregations strictly to past sliding windows reduced PR-AUC to an honest $0.0354$. We report this transparently to caution fellow practitioners against non-causal graph aggregations.

\textbf{3. Treating Absence as Information in Biometrics:} When direct script attacks bypass the browser, keystroke events are missing entirely. Imputing values via mean or median imputation softens anomalous signals. Treating telemetry absence as an explicit feature ($H = 0.00$) enables fast-path short-circuiting in Layer~3 without activating heavier inference stages.

\section{Limitations \& Future Work}
\label{sec:limitations}

While RazorVigil achieves sub-50\,ms synchronous latency and strong zero-day resilience, three limitations must be recognized:
\begin{enumerate}
    \item \textbf{Edge-Case Traveler Friction:} The $10.6\%$ FPR on legitimate VPN users requires robust fallback paths (UPI QR step-up) to avoid merchant conversion degradation.
    \item \textbf{Synthetic Training Dependency:} While our benchmark models empirical Indian payment rails under strict leakage controls, live production traffic contains idiosyncratic merchant nuances.
    \item \textbf{Graph Re-Partitioning Lag:} Louvain modularity calculations are updated asynchronously every 24 hours. Fast-moving bot rings operating within a 2-hour window must be caught by Layer~2 sliding velocity before graph topological updates mature.
\end{enumerate}
Future work will examine online streaming graph clustering via distributed message queues and evaluate federated privacy-preserving model updates across participating payment aggregators.

\section{Conclusion}
\label{sec:conclusion}

In this work, we presented \textbf{RazorVigil}, a multi-modal, nine-layer synchronous fraud interception framework engineered for Indian digital commerce under a strict 50\,ms latency SLA. By introducing a \textbf{persistence-gated ensemble} that requires unsupervised anomaly consensus in ambiguous prediction regions, RazorVigil achieves $76.8\%$ recall on unobserved zero-day CVV cycling attacks---a $9.4\times$ gain over standard supervised gradient-boosted ensembles. Evaluated on an entity-disjoint benchmark with zero card or IP leakage, RazorVigil records a held-out PR-AUC of 0.9963, an adversarial bot catch rate of $97.0\%$, and a median latency of 9.48\,ms at 40\,RPS. Rigorous cold-transfer benchmarking across 402,915 external transactions from IEEE-CIS and ULB European datasets confirms the system's identity as a specialized behavioral telemetry interceptor. By pairing formal split conformal coverage guarantees with Bayesian minimum expected loss routing, RazorVigil provides a production-grade blueprint for safeguarding high-throughput real-time payment gateways.

\bibliographystyle{IEEEtran}
\bibliography{razorvigil}

\end{document}
